% ============================================================
% fig_gradient_descent.tex: Descenso del gradiente
% Uso: % ============================================================
% fig_gradient_descent.tex: Descenso del gradiente
% Uso: % ============================================================
% fig_gradient_descent.tex: Descenso del gradiente
% Uso: % ============================================================
% fig_gradient_descent.tex: Descenso del gradiente
% Uso: \input{figures/fig_gradient_descent}
% ============================================================
\begin{figure}[H]
\centering
\begin{tikzpicture}[scale=1.1]

% Ejes
\draw[-{Stealth}, thick] (-0.3, 0) -- (6.5, 0) node[right] {$w$};
\draw[-{Stealth}, thick] (0, -0.3) -- (0, 4.2) node[above] {$\mathcal{L}(w)$};

% Curva de pérdida (parábola suavizada)
\draw[blue!70, very thick, domain=0.3:6.1, samples=80]
    plot (\x, {0.35*(\x-3.2)^2 + 0.5});

% Mínimo
\fill[green!60!black] (3.2, 0.5) circle (3pt) node[below right, font=\small] {mínimo};

% Puntos del descenso
\coordinate (p1) at (1.0, 2.34);
\coordinate (p2) at (1.7, 1.44);
\coordinate (p3) at (2.3, 0.87);
\coordinate (p4) at (2.75, 0.59);
\coordinate (p5) at (3.05, 0.51);

\fill[red!70]    (p1) circle (3pt);
\fill[red!60]    (p2) circle (3pt);
\fill[orange!80] (p3) circle (3pt);
\fill[orange!60] (p4) circle (3pt);
\fill[yellow!80!orange] (p5) circle (3pt);

% Flechas de actualización
\draw[-{Stealth}, red!60, thick]    (p1) -- (p2);
\draw[-{Stealth}, red!40, thick]    (p2) -- (p3);
\draw[-{Stealth}, orange!80, thick] (p3) -- (p4);
\draw[-{Stealth}, orange!50, thick] (p4) -- (p5);

% Etiqueta tasa de aprendizaje
\draw[dashed, gray] (1.0, 0) -- (p1);
\draw[dashed, gray] (1.7, 0) -- (p2);
\node[font=\tiny, gray] at (1.35, -0.2) {$\eta \cdot \nabla\mathcal{L}$};
\draw[{Stealth}-{Stealth}, gray, thin] (1.05, -0.15) -- (1.65, -0.15);

% Fórmula
\node[draw=blue!30, fill=blue!5, rounded corners, font=\small,
      inner sep=6pt] at (5.0, 3.5)
      {$w \leftarrow w - \eta \dfrac{\partial \mathcal{L}}{\partial w}$};

\node[font=\footnotesize, gray] at (3.2, -0.5) {$w^*$};

\end{tikzpicture}
\caption{Descenso del gradiente. Partiendo de un valor inicial de $w$, cada
iteración actualiza el parámetro en dirección opuesta al gradiente de la
pérdida $\mathcal{L}$, con magnitud controlada por la tasa de aprendizaje
$\eta$. El proceso converge progresivamente al mínimo $w^*$.}
\label{fig:gradient_descent}
\end{figure}

% ============================================================
\begin{figure}[H]
\centering
\begin{tikzpicture}[scale=1.1]

% Ejes
\draw[-{Stealth}, thick] (-0.3, 0) -- (6.5, 0) node[right] {$w$};
\draw[-{Stealth}, thick] (0, -0.3) -- (0, 4.2) node[above] {$\mathcal{L}(w)$};

% Curva de pérdida (parábola suavizada)
\draw[blue!70, very thick, domain=0.3:6.1, samples=80]
    plot (\x, {0.35*(\x-3.2)^2 + 0.5});

% Mínimo
\fill[green!60!black] (3.2, 0.5) circle (3pt) node[below right, font=\small] {mínimo};

% Puntos del descenso
\coordinate (p1) at (1.0, 2.34);
\coordinate (p2) at (1.7, 1.44);
\coordinate (p3) at (2.3, 0.87);
\coordinate (p4) at (2.75, 0.59);
\coordinate (p5) at (3.05, 0.51);

\fill[red!70]    (p1) circle (3pt);
\fill[red!60]    (p2) circle (3pt);
\fill[orange!80] (p3) circle (3pt);
\fill[orange!60] (p4) circle (3pt);
\fill[yellow!80!orange] (p5) circle (3pt);

% Flechas de actualización
\draw[-{Stealth}, red!60, thick]    (p1) -- (p2);
\draw[-{Stealth}, red!40, thick]    (p2) -- (p3);
\draw[-{Stealth}, orange!80, thick] (p3) -- (p4);
\draw[-{Stealth}, orange!50, thick] (p4) -- (p5);

% Etiqueta tasa de aprendizaje
\draw[dashed, gray] (1.0, 0) -- (p1);
\draw[dashed, gray] (1.7, 0) -- (p2);
\node[font=\tiny, gray] at (1.35, -0.2) {$\eta \cdot \nabla\mathcal{L}$};
\draw[{Stealth}-{Stealth}, gray, thin] (1.05, -0.15) -- (1.65, -0.15);

% Fórmula
\node[draw=blue!30, fill=blue!5, rounded corners, font=\small,
      inner sep=6pt] at (5.0, 3.5)
      {$w \leftarrow w - \eta \dfrac{\partial \mathcal{L}}{\partial w}$};

\node[font=\footnotesize, gray] at (3.2, -0.5) {$w^*$};

\end{tikzpicture}
\caption{Descenso del gradiente. Partiendo de un valor inicial de $w$, cada
iteración actualiza el parámetro en dirección opuesta al gradiente de la
pérdida $\mathcal{L}$, con magnitud controlada por la tasa de aprendizaje
$\eta$. El proceso converge progresivamente al mínimo $w^*$.}
\label{fig:gradient_descent}
\end{figure}

% ============================================================
\begin{figure}[H]
\centering
\begin{tikzpicture}[scale=1.1]

% Ejes
\draw[-{Stealth}, thick] (-0.3, 0) -- (6.5, 0) node[right] {$w$};
\draw[-{Stealth}, thick] (0, -0.3) -- (0, 4.2) node[above] {$\mathcal{L}(w)$};

% Curva de pérdida (parábola suavizada)
\draw[blue!70, very thick, domain=0.3:6.1, samples=80]
    plot (\x, {0.35*(\x-3.2)^2 + 0.5});

% Mínimo
\fill[green!60!black] (3.2, 0.5) circle (3pt) node[below right, font=\small] {mínimo};

% Puntos del descenso
\coordinate (p1) at (1.0, 2.34);
\coordinate (p2) at (1.7, 1.44);
\coordinate (p3) at (2.3, 0.87);
\coordinate (p4) at (2.75, 0.59);
\coordinate (p5) at (3.05, 0.51);

\fill[red!70]    (p1) circle (3pt);
\fill[red!60]    (p2) circle (3pt);
\fill[orange!80] (p3) circle (3pt);
\fill[orange!60] (p4) circle (3pt);
\fill[yellow!80!orange] (p5) circle (3pt);

% Flechas de actualización
\draw[-{Stealth}, red!60, thick]    (p1) -- (p2);
\draw[-{Stealth}, red!40, thick]    (p2) -- (p3);
\draw[-{Stealth}, orange!80, thick] (p3) -- (p4);
\draw[-{Stealth}, orange!50, thick] (p4) -- (p5);

% Etiqueta tasa de aprendizaje
\draw[dashed, gray] (1.0, 0) -- (p1);
\draw[dashed, gray] (1.7, 0) -- (p2);
\node[font=\tiny, gray] at (1.35, -0.2) {$\eta \cdot \nabla\mathcal{L}$};
\draw[{Stealth}-{Stealth}, gray, thin] (1.05, -0.15) -- (1.65, -0.15);

% Fórmula
\node[draw=blue!30, fill=blue!5, rounded corners, font=\small,
      inner sep=6pt] at (5.0, 3.5)
      {$w \leftarrow w - \eta \dfrac{\partial \mathcal{L}}{\partial w}$};

\node[font=\footnotesize, gray] at (3.2, -0.5) {$w^*$};

\end{tikzpicture}
\caption{Descenso del gradiente. Partiendo de un valor inicial de $w$, cada
iteración actualiza el parámetro en dirección opuesta al gradiente de la
pérdida $\mathcal{L}$, con magnitud controlada por la tasa de aprendizaje
$\eta$. El proceso converge progresivamente al mínimo $w^*$.}
\label{fig:gradient_descent}
\end{figure}

% ============================================================
\begin{figure}[H]
\centering
\begin{tikzpicture}[scale=1.1]

% Ejes
\draw[-{Stealth}, thick] (-0.3, 0) -- (6.5, 0) node[right] {$w$};
\draw[-{Stealth}, thick] (0, -0.3) -- (0, 4.2) node[above] {$\mathcal{L}(w)$};

% Curva de pérdida (parábola suavizada)
\draw[blue!70, very thick, domain=0.3:6.1, samples=80]
    plot (\x, {0.35*(\x-3.2)^2 + 0.5});

% Mínimo
\fill[green!60!black] (3.2, 0.5) circle (3pt) node[below right, font=\small] {mínimo};

% Puntos del descenso
\coordinate (p1) at (1.0, 2.34);
\coordinate (p2) at (1.7, 1.44);
\coordinate (p3) at (2.3, 0.87);
\coordinate (p4) at (2.75, 0.59);
\coordinate (p5) at (3.05, 0.51);

\fill[red!70]    (p1) circle (3pt);
\fill[red!60]    (p2) circle (3pt);
\fill[orange!80] (p3) circle (3pt);
\fill[orange!60] (p4) circle (3pt);
\fill[yellow!80!orange] (p5) circle (3pt);

% Flechas de actualización
\draw[-{Stealth}, red!60, thick]    (p1) -- (p2);
\draw[-{Stealth}, red!40, thick]    (p2) -- (p3);
\draw[-{Stealth}, orange!80, thick] (p3) -- (p4);
\draw[-{Stealth}, orange!50, thick] (p4) -- (p5);

% Etiqueta tasa de aprendizaje
\draw[dashed, gray] (1.0, 0) -- (p1);
\draw[dashed, gray] (1.7, 0) -- (p2);
\node[font=\tiny, gray] at (1.35, -0.2) {$\eta \cdot \nabla\mathcal{L}$};
\draw[{Stealth}-{Stealth}, gray, thin] (1.05, -0.15) -- (1.65, -0.15);

% Fórmula
\node[draw=blue!30, fill=blue!5, rounded corners, font=\small,
      inner sep=6pt] at (5.0, 3.5)
      {$w \leftarrow w - \eta \dfrac{\partial \mathcal{L}}{\partial w}$};

\node[font=\footnotesize, gray] at (3.2, -0.5) {$w^*$};

\end{tikzpicture}
\caption{Descenso del gradiente. Partiendo de un valor inicial de $w$, cada
iteración actualiza el parámetro en dirección opuesta al gradiente de la
pérdida $\mathcal{L}$, con magnitud controlada por la tasa de aprendizaje
$\eta$. El proceso converge progresivamente al mínimo $w^*$.}
\label{fig:gradient_descent}
\end{figure}
